% Source: https://artofproblemsolving.com/community/c3h2357585_666th_post_special__epic_math_problems
% Stable identity: aops-topic-2357585-q06
Join the centre $O$ to the eight vertices. This cuts the octagon into $8$ congruent isosceles triangles with two sides $5$ and apex angle $\dfrac{360^\circ}{8}=45^\circ$.

\begin{asy}
usepackage("amsmath");
size(5cm);
defaultpen(fontsize(10pt));
pair O=(0,0);
draw(circle(O,5),gray);
pair[] V; for(int k=0;k<8;++k) V.push(5*dir(22.5+45*k));
fill(O--V[0]--V[1]--cycle,paleblue);
draw(V[0]--V[1]--V[2]--V[3]--V[4]--V[5]--V[6]--V[7]--cycle,linewidth(0.9));
for(int k=0;k<8;++k) draw(O--V[k],dotted);
draw(O--V[0],linewidth(0.8)); draw(O--V[1],linewidth(0.8));
pair Fo=dot(V[1],unit(V[0]))*unit(V[0]);
draw(V[1]--Fo,red+dashed);
dot(O); label("$O$",O,W);
label("$5$",(O+V[0])/2,S); label("$5$",(O+V[1])/2,NW);
draw(arc(O,1.1,22.5,67.5)); label("$45^\circ$",O+1.7*dir(45),fontsize(8pt));
label("$h$",(V[1]+Fo)/2,E,red);
\end{asy}

In one of them, take one radius as the base. The height $h$ from the opposite vertex is a leg of a $45^\circ$-$45^\circ$-$90^\circ$ triangle with hypotenuse $5$, so $h=\dfrac{5}{\sqrt2}=\dfrac{5\sqrt2}{2}$. Each triangle has area $\dfrac12\cdot5\cdot\dfrac{5\sqrt2}{2}=\dfrac{25\sqrt2}{4}$, and the octagon has area
\[8\cdot\dfrac{25\sqrt2}{4}=\boxed{50\sqrt2}.\]
